\documentclass[10pt,a4paper]{article}
\usepackage[T1]{fontenc}
\usepackage[utf8]{inputenc}
\usepackage[english]{babel}
\usepackage{lmodern}
\usepackage{amsmath,amssymb}
\usepackage{xcolor}
\usepackage[left=2.25cm,right=2.25cm,top=1.65cm,bottom=1.65cm]{geometry}
\usepackage{microtype}
\definecolor{PFFBlue}{HTML}{203E6A}
\setlength{\parindent}{0pt}
\setlength{\parskip}{0.23em}
\setlength{\abovedisplayskip}{6pt}
\setlength{\belowdisplayskip}{6pt}
\begin{document}

\section*{\textcolor{PFFBlue}{HW 44 --- Properties of Integrals of Simple Functions}}

\textbf{Statement.} Fix a measure space $(S,\Sigma,\mu)$. Prove the following properties.

\textbf{(i)} If $f,g\in SF^+$ and $\mu(\{s\in S:f(s)\ne g(s)\})=0$, then $\mu_0(f)=\mu_0(g)$.

\textbf{(ii)} If $f,g\in SF^+$ and $c\ge0$, then $f+g,cf\in SF^+$; moreover,
\[
\mu_0(f+g)=\mu_0(f)+\mu_0(g),\qquad \mu_0(cf)=c\mu_0(f).
\]

\textbf{(iii)} If $f,g\in SF^+$ and $f\le g$, then $\mu_0(f)\le\mu_0(g)$.

\textbf{(iv)} If $f,g\in SF^+$, then $f\vee g:=\max\{f,g\}$ and $f\wedge g:=\min\{f,g\}$ also belong to $SF^+$.

\medskip
\textbf{Recall --- $\boldsymbol{\mu}$-almost everywhere.}
A property $P(s)$ holds \emph{$\mu$-almost everywhere} ($\mu$-a.e.) if it fails only on a set of $\mu$-measure zero. In particular, for measurable $f,g$,
\[
f=g\quad\mu\text{-a.e.}
\quad\Longleftrightarrow\quad
\mu\bigl(\{s\in S:f(s)\ne g(s)\}\bigr)=0.
\]
Thus $f$ and $g$ need not agree at every point.

When constructing a common partition as in HW~43, we discard only \emph{empty} intersections, since partition elements must be nonempty. We do not discard nonempty sets of measure zero: for example, on $S=[0,1]$ with $\mu=\mathrm{Leb}$,
\[
D=\{\tfrac12\}\ne\emptyset,\qquad \mu(D)=0.
\]
Such a set remains a partition element, although its contribution to the integral is zero.

\medskip
\textbf{\textcolor{PFFBlue}{Proof.}}

By HW~43, both $f$ and $g$ admit canonical representations, whose level sets form finite measurable partitions of $S$. Intersecting these partitions and omitting empty intersections gives a common finite measurable partition $D_1,\ldots,D_r$, on which both functions are constant.
\[
f=\sum_{k=1}^{r}a_k\mathbf1_{D_k},\qquad
 g=\sum_{k=1}^{r}b_k\mathbf1_{D_k},
\]
where $a_k,b_k\ge0$. By HW~43, the integral is independent of the representation, so
\[
\mu_0(f)=\sum_{k=1}^{r}a_k\mu(D_k),\qquad
\mu_0(g)=\sum_{k=1}^{r}b_k\mu(D_k).
\]

\textbf{(i) Equality almost everywhere.}
Suppose $f=g$ $\mu$-a.e., and let
\[
N:=\{s\in S:f(s)\ne g(s)\},\qquad \mu(N)=0.
\]
Fix $k$ such that $a_k\ne b_k$. For any $s\in D_k$, the partition representations give
\[
f(s)=a_k\ne b_k=g(s)\quad\Longrightarrow\quad s\in N.
\]
Since $s\in D_k$ was arbitrary, $D_k\subseteq N$. By monotonicity of $\mu$,
\[
0\le\mu(D_k)\le\mu(N)=0
\quad\Longrightarrow\quad \mu(D_k)=0.
\]
Consequently, for every $k=1,\ldots,r$,
\[
a_k\mu(D_k)=b_k\mu(D_k).
\]
Indeed, if $a_k=b_k$ the equality is immediate; otherwise $\mu(D_k)=0$ and both sides vanish. Summing gives
\[
\mu_0(f)=\sum_{k=1}^{r}a_k\mu(D_k)
=\sum_{k=1}^{r}b_k\mu(D_k)=\mu_0(g).
\]

\newpage
\textbf{(ii) Additivity and positive homogeneity.}
Using the common partition,
\[
f+g=\sum_{k=1}^{r}(a_k+b_k)\mathbf1_{D_k},\qquad
cf=\sum_{k=1}^{r}(ca_k)\mathbf1_{D_k}.
\]
Thus $f+g,cf\in SF^+$, since their coefficients are nonnegative and the $D_k$ are measurable. Moreover,
\[
\begin{aligned}
\mu_0(f+g)
&=\sum_{k=1}^{r}(a_k+b_k)\mu(D_k)\\
&=\sum_{k=1}^{r}a_k\mu(D_k)
 +\sum_{k=1}^{r}b_k\mu(D_k)\\
&=\mu_0(f)+\mu_0(g),
\end{aligned}
\]
and
\[
\begin{aligned}
\mu_0(cf)
&=\sum_{k=1}^{r}(ca_k)\mu(D_k)\\
&=c\sum_{k=1}^{r}a_k\mu(D_k)
=c\mu_0(f).
\end{aligned}
\]
The equalities follow from the definition of $\mu_0$ and distributivity of finite nonnegative sums, with $0\cdot\infty=0$ when necessary.

\textbf{(iii) Monotonicity.}
Suppose $f\le g$. Since $D_k\ne\emptyset$ and both functions are constant on it, $a_k\le b_k$ for every $k=1,\ldots,r$. Hence
\[
\mu_0(f)=\sum_{k=1}^{r}a_k\mu(D_k)
\le\sum_{k=1}^{r}b_k\mu(D_k)=\mu_0(g).
\]
The inequality follows from $a_k\le b_k$ and $\mu(D_k)\ge0$ for every $k$.

\textbf{(iv) Maximum and minimum.}
Since $f=a_k$ and $g=b_k$ on $D_k$,
\[
\begin{aligned}
f\vee g&=\sum_{k=1}^{r}\max\{a_k,b_k\}\mathbf1_{D_k},\\
f\wedge g&=\sum_{k=1}^{r}\min\{a_k,b_k\}\mathbf1_{D_k}.
\end{aligned}
\]
Both functions belong to $SF^+$, because their coefficients are nonnegative and their partition elements are measurable.

Therefore, the statement is proved. $\square$
\end{document}
